\documentclass[11pt,a4paper]{article}
\usepackage[utf8]{inputenc}
\usepackage[T1]{fontenc}
\usepackage[portuguese]{babel}
\usepackage{geometry}
\geometry{margin=2.5cm}
\usepackage{listings}
\usepackage{xcolor}
\usepackage{amsmath}
\usepackage{amssymb}
\usepackage{hyperref}
\usepackage{enumitem}
\usepackage{tcolorbox}
\usepackage{tabularx}
\usepackage{booktabs}

\definecolor{codebg}{rgb}{0.95,0.95,0.95}
\definecolor{codegreen}{rgb}{0.1,0.5,0.1}
\definecolor{codepurple}{rgb}{0.5,0.1,0.5}
\definecolor{codegray}{rgb}{0.5,0.5,0.5}

\lstdefinestyle{python}{
    backgroundcolor=\color{codebg},
    commentstyle=\color{codegray}\itshape,
    keywordstyle=\color{codepurple}\bfseries,
    stringstyle=\color{codegreen},
    basicstyle=\ttfamily\small,
    breaklines=true,
    frame=single,
    language=Python,
    showstringspaces=false,
    tabsize=4,
    literate={á}{{\'a}}1 {ã}{{\~a}}1 {é}{{\'e}}1 {í}{{\'i}}1 {ó}{{\'o}}1 {õ}{{\~o}}1 {ú}{{\'u}}1 {ç}{{\c{c}}}1 {Á}{{\'A}}1 {Ã}{{\~A}}1 {É}{{\'E}}1 {Í}{{\'I}}1 {Ó}{{\'O}}1 {Õ}{{\~O}}1 {Ú}{{\'U}}1 {Ç}{{\c{C}}}1 {â}{{\^a}}1 {ê}{{\^e}}1 {ô}{{\^o}}1 {û}{{\^u}}1 {Â}{{\^A}}1 {Ê}{{\^E}}1 {Ô}{{\^O}}1 {Û}{{\^U}}1 {à}{{\`a}}1 {è}{{\`e}}1 {ì}{{\`i}}1 {ò}{{\`o}}1 {ù}{{\`u}}1
}
\lstset{style=python}

\newtcolorbox{solucao}{
  colback=yellow!5!white,
  colframe=orange!60!black,
  title=Solução proposta,
  fonttitle=\bfseries
}

\title{\textbf{Data Analysis with Python 2024--2025}\\Parte 3: Mais NumPy\\\large Soluções dos Exercícios}
\author{Soluções independentes baseadas nos enunciados originais}
\date{}

\begin{document}
\maketitle

\section*{Nota Importante}
Este material é uma adaptação das soluções independentes para os exercícios baseados no curso \textit{Data Analysis with Python 2024-2025}, oferecido pela \textbf{The University of Helsinki MOOC Center}. As soluções abaixo não são o gabarito oficial do curso.

\tableofcontents
\newpage

\section{Soluções - Capítulo 3}

\subsection*{Exercício 3.11 -- para tons de cinza (\textit{to grayscale})}

\begin{solucao}
\begin{lstlisting}
import numpy as np
import matplotlib.pyplot as plt

def to_grayscale(image):
    # Pesos definidos no problema: r=0.2126, g=0.7152, b=0.0722
    weights = np.array([0.2126, 0.7152, 0.0722])
    # Produto escalar para obter a matriz 2D final em tons de cinza
    return np.dot(image[..., :3], weights)

def to_red(image):
    img = image.copy()
    img[:, :, 1] = 0  # Zera o verde
    img[:, :, 2] = 0  # Zera o azul
    return img

def to_green(image):
    img = image.copy()
    img[:, :, 0] = 0  # Zera o vermelho
    img[:, :, 2] = 0  # Zera o azul
    return img

def to_blue(image):
    img = image.copy()
    img[:, :, 0] = 0  # Zera o vermelho
    img[:, :, 1] = 0  # Zera o verde
    return img

def main():
    painting = plt.imread('src/painting.png')
    
    gray = to_grayscale(painting)
    plt.gray()
    plt.imshow(gray)
    plt.show()

    fig, ax = plt.subplots(3, 1)
    ax[0].imshow(to_red(painting))
    ax[1].imshow(to_green(painting))
    ax[2].imshow(to_blue(painting))
    plt.show()

if __name__ == '__main__':
    main()
\end{lstlisting}
\end{solucao}

\subsection*{Exercício 3.12 -- desvanecimento radial (\textit{radial fade})}

\begin{solucao}
\begin{lstlisting}
import numpy as np
import matplotlib.pyplot as plt

def center(a):
    return (a.shape[0] - 1) / 2.0, (a.shape[1] - 1) / 2.0

def radial_distance(a):
    h, w = a.shape[:2]
    cy, cx = center(a)
    y, x = np.ogrid[0:h, 0:w]
    return np.sqrt((x - cx)**2 + (y - cy)**2)

def scale(a, tmin=0.0, tmax=1.0):
    a_min = np.min(a)
    a_max = np.max(a)
    if a_min == a_max:
        return np.full_like(a, tmin)
    return ((a - a_min) / (a_max - a_min)) * (tmax - tmin) + tmin

def radial_mask(a):
    dist = radial_distance(a)
    scaled_dist = scale(dist)
    return 1 - scaled_dist

def radial_fade(a):
    mask = radial_mask(a)
    # Garante o broadcasting correto com a terceira dimensao (cores)
    return a * mask[..., np.newaxis]

def main():
    painting = plt.imread('src/painting.png')
    
    fig, ax = plt.subplots(3, 1)
    ax[0].imshow(painting)
    ax[1].imshow(radial_mask(painting), cmap='gray')
    ax[2].imshow(radial_fade(painting))
    plt.show()

if __name__ == '__main__':
    main()
\end{lstlisting}
\end{solucao}

\vspace{2em}
\noindent\textbf{Créditos:} Este conteúdo foi originalmente criado pela \textit{The University of Helsinki MOOC Center} para o curso \textit{Data Analysis with Python}.
\end{document}
